\documentclass[
  reprint,
  superscriptaddress,
  amsmath,
  amssymb,
  aps,
  pre,
  longbibliography
]{revtex4-2}

% Keep \bm lighter by reserving fewer bold math alphabets.
\newcommand\hmmax{0}
\newcommand\bmmax{0}

% Text and math fonts.
\usepackage[utf8]{inputenc}
\usepackage[lining,semibold]{libertine}
\usepackage{amsthm}
\newtheorem{theorem}{Theorem}
\newtheorem{proposition}[theorem]{Proposition}
\usepackage[libertine,cmintegrals,bigdelims,vvarbb]{newtxmath}

% Core math support.
\usepackage{amsmath}
\usepackage{amsfonts}
\usepackage{bm}
\usepackage{needspace}

% Figures and tables.
\usepackage{graphicx}
% REVTeX's table patches require the legacy array implementation on TeX Live 2025.
% Load it before dcolumn, which otherwise loads the incompatible current array.
\usepackage{array}[=2016-10-06]
\usepackage{dcolumn}
\usepackage{booktabs}
\usepackage{siunitx}

% Links and references. Load hyperref before cleveref.
\usepackage{xcolor}
\definecolor{webgreen}{rgb}{0,.5,0}
\definecolor{webbrown}{rgb}{.6,0,0}
\definecolor{RoyalBlue}{rgb}{0.0,0.14,0.4}

\usepackage{hyperref}
\hypersetup{
  colorlinks=true,
  linktocpage=true,
  breaklinks=true,
  pdfstartpage=1,
  pdfstartview=Fit,
  pdfpagemode=UseOutlines,
  plainpages=false,
  bookmarksnumbered=true,
  bookmarksopen=true,
  bookmarksopenlevel=1,
  hypertexnames=true,
  pdfhighlight=/O,
  urlcolor=webbrown,
  linkcolor=RoyalBlue,
  filecolor=RoyalBlue,
  citecolor=webgreen,
  pdftitle={Mutual linearity from classical Markov response theory},
  pdfauthor={Timur Aslyamov and Massimiliano Esposito},
  pdfcreator={pdfLaTeX},
  pdfproducer={LaTeX REVTeX}
}

\usepackage[capitalize]{cleveref}
\newcommand{\creflastconjunction}{, and\nobreakspace}
\newcommand{\crefrangeconjunction}{--}

% Prevent math macros from leaking into PDF bookmarks.
\pdfstringdefDisableCommands{%
  \def\theta{}%
  \def\varphi{}%
  \def\mathcal{}%
  \def\mathbb{}%
}

% Normal-size equation numbers.
\makeatletter
\def\maketag@@@#1{\hbox{\m@th\normalfont\normalsize#1}}
\makeatother

% Small reusable macros.
\newcommand{\me}{\mathrm{e}}
\newcommand{\tr}{\operatorname{tr}}
\newcommand{\mb}[1]{\mathbf{#1}}
\newcommand{\E}{\mathbb{E}}
\newcommand{\Prob}{\mathbb{P}}
\newcommand{\dd}{\mathrm{d}}
\newcommand{\src}{\operatorname{s}}
\newcommand{\tgt}{\operatorname{t}}


\begin{document}
%%%%%%%%%%%%%%%%%%%%%%%%%%%%%%%%%%%%% frontmatter %%%%%%%%%%%%%%%%%%%%%%%%%%%%%%%%%%%%%
\title{Mutual linearity from classical Markov response theory}

\author{Timur Aslyamov}
\email{timur.aslyamov@uni.lu}
\affiliation{Complex Systems and Statistical Mechanics, Department of Physics and Materials Science, University of Luxembourg, 30 Avenue des Hauts-Fourneaux, L-4362 Esch-sur-Alzette, Luxembourg}

\author{Massimiliano Esposito}
\email{massimiliano.esposito@uni.lu}
\affiliation{Complex Systems and Statistical Mechanics, Department of Physics and Materials Science, University of Luxembourg, 30 Avenue des Hauts-Fourneaux, L-4362 Esch-sur-Alzette, Luxembourg}

\date{\today}


\begin{abstract}
The exact change of a stationary Markov distribution under a finite
perturbation is a classical result of Markov response theory. Written with
the Drazin inverse, it shows immediately that varying the two rates of
one edge confines the stationary distribution to a fixed affine line.
We use this observation to derive mutual linearity of probabilities,
fixed linear observables, and stationary currents, and to identify their
coefficients with integrated relaxation and first-passage times.
More generally, a perturbation whose range lies in a fixed
$r$-dimensional subspace confines stationary distributions to an affine
subspace of dimension at most $r$. This provides a common algebraic
basis for recent mutual-linearity results and makes their assumptions
and degeneracies explicit.
\end{abstract}
\maketitle

\textit{Exact finite response.---}The starting point is the classical
finite-perturbation identity of Markov-chain theory
\cite{Schweitzer1968,Meyer1975,Meyer1980}. Consider two irreducible
continuous-time Markov jump processes on the same finite state space.
We use column probability vectors, $\dot{\bm p}=W\bm p$, so that
$W_{ij}$ is the rate from $j$ to $i$ for $i\ne j$ and
$\bm 1^{\mathsf T}W=0$. Their normalized stationary distributions
obey $W\bm\pi=0$ and $W'\bm\pi'=0$. For an arbitrary finite
change $\Delta W=W'-W$,
\begin{equation}
 \boxed{\bm\pi'-\bm\pi=-W^D\Delta W\bm\pi'.}
 \label{eq:finite}
\end{equation}
Here $G\equiv W^D$ is the Drazin inverse, which for an irreducible
finite generator equals the group inverse. With
$\Pi=\bm\pi\bm 1^{\mathsf T}$, it is characterized by
\begin{equation}
 WG=GW=I-\Pi,\qquad G\bm\pi=0,\qquad
 \bm 1^{\mathsf T}G=0.
 \label{eq:group}
\end{equation}
Indeed, subtracting stationarity equations gives
$W(\bm\pi'-\bm\pi)=-\Delta W\bm\pi'$. Applying $G$ and using
$\bm 1^{\mathsf T}(\bm\pi'-\bm\pi)=0$ proves
\cref{eq:finite}. There is no expansion in $\Delta W$.

To connect directly with the fundamental matrix, choose a rate
$\nu$ larger than every escape rate of both processes and set
$T=I+W/\nu$, $T'=I+W'/\nu$. Then
\begin{equation}
 Z=(I-T+\Pi)^{-1}=\Pi-\nu G,
 \qquad
 \bm\pi'-\bm\pi=Z(T'-T)\bm\pi'.
 \label{eq:fundamental}
\end{equation}
This is the discrete-time identity in column convention
\cite{Schweitzer1968,Meyer1980}; rank-one stationary updates were also
studied explicitly by Hunter~\cite{Hunter1986}.
An ordinary Moore--Penrose inverse $W^+$ can be used, but its solution
must be normalized:
$\bm\pi'-\bm\pi=-(I-\Pi)W^+\Delta W\bm\pi'$.
In general $W^+$ and $W^D$ are different.

Recent work established affine relations between stationary network
currents~\cite{Harunari2024}, interpreted their coefficients through
first-passage times~\cite{Khodabandehlou2025}, and extended them to
probabilities and other observables~\cite{BebonSpeck2026} and to multiple
controlled edges~\cite{DalCengio2025}. We show that their stationary
linearity relations follow from \cref{eq:finite} and the range of the
generator perturbation. The purpose is to make this classical basis
and its consequences explicit.

\textit{One edge, one response direction.---}Orient the controlled edge
$\ell$ from state $a$ to state $b$. Write
\begin{equation}
 \bm d=\bm e_b-\bm e_a,\qquad
 \bm c=k_+\bm e_a-k_-\bm e_b,
 \qquad j_\ell=\bm c^{\mathsf T}\bm\pi,
 \label{eq:edge}
\end{equation}
where $\bm e_a$ is a coordinate vector. The edge contributes
$\bm d\bm c^{\mathsf T}$ to $W$, including its diagonal entries.
Changing either or both rates, with all other rates fixed, gives
\begin{equation}
 \Delta W=\bm d\bm v^{\mathsf T},\qquad
 \bm v=\Delta k_+\bm e_a-\Delta k_-\bm e_b.
 \label{eq:rankone}
\end{equation}
Both rate changes have the same range, $\operatorname{span}\{\bm d\}$.
Thus \cref{eq:finite} reduces to
\begin{equation}
 \boxed{\bm\pi'=\bm\pi+\alpha\bm h,\qquad
 \bm h=-G\bm d,\qquad \alpha=\bm v^{\mathsf T}\bm\pi'.}
 \label{eq:line}
\end{equation}
The vector $\bm h$ is fixed by any chosen reference generator;
all dependence on the changed rates is in one scalar $\alpha$.
In particular,
\begin{equation}
 \alpha=\frac{\bm v^{\mathsf T}\bm\pi}
                  {1-\bm v^{\mathsf T}\bm h}.
 \label{eq:amplitude}
\end{equation}
The denominator is nonzero whenever $W'$ remains irreducible:
if it vanished, the nonzero zero-sum vector $\bm h$ would satisfy
$W'\bm h=-\bm d+\bm d\bm v^{\mathsf T}\bm h=0$, contradicting
uniqueness of the normalized stationary state.
The nonlinear rate dependence in \cref{eq:amplitude} is fully
compatible with the affine geometry in \cref{eq:line}.

\textit{Probabilities and observables.---}For any states $m,n$,
\begin{equation}
 h_m(\pi_n'-\pi_n)=h_n(\pi_m'-\pi_m).
 \label{eq:probabilities}
\end{equation}
When $h_m\ne0$, this gives the probability--probability slope
$h_n/h_m$ and offset $\pi_n-(h_n/h_m)\pi_m$, independent of the
two controlled rates. A state with $h_m=0$ is insensitive to these
changes and cannot serve as a reference coordinate, even though
$\pi_m>0$. More generally, for fixed coefficient vectors $\bm f$
and $\bm g$, let $F=\bm f^{\mathsf T}\bm\pi$ and
$H=\bm g^{\mathsf T}\bm\pi$. Then
\begin{equation}
 F'-F=\frac{\bm f^{\mathsf T}\bm h}
                 {\bm g^{\mathsf T}\bm h}(H'-H),
 \qquad \bm g^{\mathsf T}\bm h\ne0.
 \label{eq:observables}
\end{equation}
This includes state averages, mean occupation-time rates, and mean
jump-counting rates on unchanged edges: for fixed jump weights
$q_{ij}$, $\sum_{i\ne j}q_{ij}W_{ij}\pi_j$ is a fixed linear
functional of $\bm\pi$ if the counted rates are unchanged.
These are the stationary observable relations of
Ref.~\cite{BebonSpeck2026}.

The ratio of the two rate responses follows just as directly.
Taking an infinitesimal change about any member of the family gives
\begin{equation}
 \partial_{k_+}\pi_n=h_n\pi_a,\qquad
 \partial_{k_-}\pi_n=-h_n\pi_b,
 \qquad
 -\frac{\partial_{k_+}\pi_n}{\partial_{k_-}\pi_n}
 =\frac{\pi_a}{\pi_b},
 \label{eq:response-ratio}
\end{equation}
where $G$ and $\bm h$ are evaluated at that member and the ratio
requires $h_n\ne0$. This recovers the response-ratio identity of
Ref.~\cite{BebonSpeck2026} without a spanning-tree calculation.

\textit{Stationary currents.---}For an unchanged edge $e$, write
$j_e=\bm c_e^{\mathsf T}\bm\pi$ with its fixed current vector
$\bm c_e$. The controlled current has a changing coefficient vector,
so its variation contains a direct contribution:
\begin{align}
 j_e'-j_e&=\alpha\bm c_e^{\mathsf T}\bm h,
 \label{eq:output-current}\\
 j_\ell'-j_\ell
 &=\bm c^{\mathsf T}(\bm\pi'-\bm\pi)
   +\bm v^{\mathsf T}\bm\pi'
 =\alpha\kappa,
 \qquad \kappa=1+\bm c^{\mathsf T}\bm h.
 \label{eq:input-current}
\end{align}
Eliminating $\alpha$, for $\kappa\ne0$,
\begin{equation}
 \boxed{j_e'-j_e=\chi_{e\ell}(j_\ell'-j_\ell),\qquad
 \chi_{e\ell}=
 -\frac{\bm c_e^{\mathsf T}G\bm d}
        {1-\bm c^{\mathsf T}G\bm d}.}
 \label{eq:currents}
\end{equation}
This is the mutual affine linearity of stationary currents
\cite{Harunari2024}. It holds for arbitrary simultaneous changes of
$k_+$ and $k_-$ that preserve irreducibility. The same slope and
offset $j_e-\chi_{e\ell}j_\ell$ apply throughout the family.

Their independence of the input rates is especially transparent when
the generator with edge $\ell$ deleted,
$W_0=W-\bm d\bm c^{\mathsf T}$, is irreducible. Denote its
stationary distribution and group inverse by $\bm\pi_0$ and $G_0$.
Applying \cref{eq:finite} with $W_0$ as reference gives directly
\begin{equation}
 \bm\pi'=\bm\pi_0-G_0\bm d\,j_\ell',\qquad
 j_e'=\bm c_e^{\mathsf T}\bm\pi_0
       -\bm c_e^{\mathsf T}G_0\bm d\,j_\ell'.
 \label{eq:deleted}
\end{equation}
Since $W_0\bm h=-\kappa\bm d$ and $\bm h$ is nonzero and zero-sum,
irreducibility of $W_0$ ensures $\kappa\ne0$. Thus
$\bm h/\kappa=-G_0\bm d$ and
$\chi_{e\ell}=-\bm c_e^{\mathsf T}G_0\bm d$.
For a network with positive rates in both directions on every edge,
this deletion condition is exactly that the input edge is not a bridge.
For general directed networks the remaining directed connectivity must
be checked. If $\kappa=0$, the input current does not change and
cannot parametrize the response. A bridge is the simplest example:
its stationary current is identically zero, although changing its rates
can redistribute probability. Equation~\eqref{eq:line} still holds.

Counting on the controlled edge requires the same care as
\cref{eq:input-current}. In particular, its traffic
$k_+\pi_a+k_-\pi_b$ is not generally affine in the probabilities
or in $j_\ell$ when both rates are freely varied. Already for two states,
$\bm\pi=(k_-,k_+)^{\mathsf T}/(k_++k_-)$ is unchanged by a
common rescaling of the rates, whereas the traffic
$2k_+k_-/(k_++k_-)$ changes. The fixed-coefficient condition in
\cref{eq:observables} is therefore essential.

\Needspace{9\baselineskip}
\textit{Relaxation and first-passage times.---}The response direction
has a direct dynamical meaning. For a finite irreducible generator,
\begin{equation}
 \begin{split}
 -G&=\int_0^\infty\!\bigl(\me^{Wt}-\Pi\bigr)\,\dd t,\\
 h_n&=\int_0^\infty\!\left[
 (\me^{Wt})_{nb}-(\me^{Wt})_{na}\right]\,\dd t.
 \end{split}
 \label{eq:relaxation}
\end{equation}
Changing one edge injects probability at one endpoint and removes it
at the other. The fixed unperturbed relaxation maps that source into
$\bm h$; its strength is $\alpha$.
Let $\tau_{x\to n}$ be the mean first-passage time to $n$ from $x$,
with $\tau_{n\to n}=0$. The group-inverse representation
\cite{Meyer1975} gives
\begin{equation}
 \tau_{x\to n}=\frac{G_{nx}-G_{nn}}{\pi_n},\qquad
 h_n=\pi_n(\tau_{a\to n}-\tau_{b\to n}).
 \label{eq:mfpt}
\end{equation}
For completeness, the first identity follows because the vector on its
right has value zero at $x=n$ and satisfies
$(W^{\mathsf T}\bm\tau)_x=-1$ for $x\ne n$, by \cref{eq:group}.
Substituting \cref{eq:mfpt} into \cref{eq:line} yields
\begin{equation}
 \pi_n'-\pi_n
 =\pi_n(\tau_{a\to n}-\tau_{b\to n})
       (\Delta k_+\pi_a'-\Delta k_-\pi_b').
 \label{eq:mfpt-response}
\end{equation}
This is the single-edge form of the finite-change relation used in
Ref.~\cite{Khodabandehlou2025}. That work already identifies its
connection to Eq.~(44) of Ref.~\cite{Harvey2023}. The fundamental
matrix, Drazin inverse, and first-passage descriptions express the same
stationary response. Using $G_0$ in \cref{eq:mfpt} likewise recovers
the current susceptibility in terms of first-passage times of the
network with the input edge deleted.

\textit{A fixed perturbation range.---}The argument is not specific to
an edge. Fix an irreducible reference $W$ and suppose a family of
irreducible generators has
\begin{equation}
 W(\bm\theta)-W=B V(\bm\theta)^{\mathsf T},\qquad
 \bm 1^{\mathsf T}B=0,\qquad \operatorname{rank}B=r,
 \label{eq:rankr}
\end{equation}
with $B$ independent of the controls. Equation~\eqref{eq:finite} implies
\begin{equation}
 \bm\pi(\bm\theta)-\bm\pi
 =-G B V(\bm\theta)^{\mathsf T}\bm\pi(\bm\theta)
 \in\operatorname{im}(GB).
 \label{eq:subspace}
\end{equation}
All stationary distributions, and hence any collection of fixed linear
observables, lie in an affine subspace of dimension at most $r$.
The range must be fixed throughout the family; a rank bound on each
individual perturbation alone does not imply this conclusion.
For one edge $r=1$, even though two rates can be controlled
independently. For several edges, $B$ can be their incidence matrix,
whose rank can be smaller than the number of controlled edges.

If deleting a set $\mathcal I$ of input edges leaves an irreducible
generator $W_0$, let $B$ collect their oriented incidence vectors and
$\bm j_{\mathcal I}$ their stationary currents. Since
$W\bm\pi=W_0\bm\pi+B\bm j_{\mathcal I}=0$,
\begin{equation}
 \boxed{\bm\pi=\bm\pi_0-G_0B\bm j_{\mathcal I},\qquad
 j_e=j_e^0-\bm c_e^{\mathsf T}G_0B\bm j_{\mathcal I},
 \quad e\notin\mathcal I,}
 \label{eq:multiple}
\end{equation}
where $j_e^0=\bm c_e^{\mathsf T}\bm\pi_0$.
This recovers the stationary mutual-multilinearity relation
\cite{DalCengio2025} with coefficients determined entirely by the
remaining network. These coefficients depend on the full deleted set;
in general they are not the coefficients obtained by perturbing each
input separately~\cite{Khodabandehlou2025}.

\Needspace{10\baselineskip}
\textit{Discussion.---}The essence of mutual linearity is the
combination of a fixed linear stationarity equation and a perturbation
with one fixed source direction. Its amplitude can depend nonlinearly
on the controlled rates, but eliminating that single amplitude leaves
affine relations between stationary outputs. Detailed balance is
unnecessary, and no small-perturbation assumption enters the argument.
The classical finite-response identity therefore supplies a common,
short derivation of the recent stationary linearity relations. Graphical formulas
and first-passage representations give complementary ways to evaluate
and interpret the coefficients. More generally, the dimension of the
fixed perturbation range bounds the dimension of stationary response,
independently of how many parameters generate that perturbation.

\bibliography{references}
\end{document}
